% =============================================================================
%  SDU University — Bachelor Diploma Project (Information Systems / 6B06101)
%  Project : Oqyrman — A Web-Based Digital Library and Reservation Platform
%            with an Integrated E-Reader and AI Reading Assistant
%  Team    : Oqyrmandar
%  Recommended compiler: XeLaTeX (Cyrillic + Kazakh support)
% =============================================================================

\documentclass[12pt,a4paper]{report}

% ---- Packages ---------------------------------------------------------------
\usepackage[T2A,T1]{fontenc}
\usepackage[utf8]{inputenc}
\usepackage[english,russian,kazakh]{babel}
\usepackage[a4paper,left=30mm,right=15mm,top=20mm,bottom=20mm]{geometry}
\usepackage{setspace}
\onehalfspacing
\usepackage{graphicx}
\usepackage{caption}
\usepackage{subcaption}
\usepackage{float}
\usepackage{enumitem}
\usepackage{titlesec}
\usepackage{tocloft}
\usepackage{listings}
\usepackage{xcolor}
\usepackage{hyperref}
\usepackage{array}
\usepackage{tabularx}
\usepackage{longtable}
\usepackage{booktabs}
\usepackage{indentfirst}

\hypersetup{colorlinks=true,linkcolor=black,citecolor=black,urlcolor=blue}

% ---- Chapter / section formatting -------------------------------------------
\titleformat{\chapter}[block]{\normalfont\Large\bfseries}{\thechapter}{1em}{\Large}
\titlespacing*{\chapter}{0pt}{-20pt}{20pt}
\titleformat{\section}{\normalfont\large\bfseries}{\thesection}{1em}{}
\titleformat{\subsection}{\normalfont\normalsize\bfseries}{\thesubsection}{1em}{}
\renewcommand{\cftchapleader}{\cftdotfill{\cftdotsep}}

% ---- Code listings ----------------------------------------------------------
\definecolor{codegray}{rgb}{0.5,0.5,0.5}
\definecolor{codebg}{rgb}{0.97,0.97,0.97}
\definecolor{codekw}{rgb}{0.10,0.30,0.70}
\lstset{
    basicstyle=\ttfamily\small,
    backgroundcolor=\color{codebg},
    keywordstyle=\color{codekw}\bfseries,
    commentstyle=\color{codegray}\itshape,
    breaklines=true,
    frame=single,
    numbers=left,
    numberstyle=\tiny\color{codegray},
    captionpos=b,
    showstringspaces=false,
    tabsize=2
}

% =============================================================================
\begin{document}
\selectlanguage{english}
\pagenumbering{roman}

% =============================================================================
%  TITLE PAGE 1  (matches SDU diploma cover layout)
% =============================================================================
\begin{titlepage}
    \begin{center}
        \vspace*{15mm}
        % \includegraphics[width=0.18\textwidth]{figures/sdu_logo.png}\\[6mm]
        \textit{[FIGURE PLACEHOLDER: SDU University logo, centred, width $\approx$ 3 cm.]}\\[10mm]
        {\large SDU University}\\[2mm]
        {\large Faculty of Engineering and Natural Sciences}\\[40mm]

        {\Huge \textbf{DIPLOMA PROJECT}}\\[12mm]
        {\Large \textbf{Oqyrman: A Web-Based Digital Library and\\
        Reservation Platform with an Integrated\\
        E-Reader and AI Reading Assistant}}\\[20mm]

        {\large Gainolla Azat, Arginbekov Alisher,\\
        Kengesbek Zanggar, Nassikhov Alikhan}\\[6mm]
        {\large 6B06101 --- ``Information Systems''}

        \vfill
        \textbf{Kaskelen, 2026}
    \end{center}
\end{titlepage}

% =============================================================================
%  TITLE PAGE 2 (Dean / Supervisor / Co-supervisor / Students)
% =============================================================================
\begin{titlepage}
    \begin{center}
        \vspace*{10mm}
        \textit{[FIGURE PLACEHOLDER: SDU University logo, centred, width $\approx$ 3 cm.]}\\[8mm]
        {\large SDU University}\\[2mm]
        {\large Faculty of Engineering and Natural Sciences}\\[20mm]
    \end{center}

    \begin{flushright}
        Dean of Faculty\\
        Assistant Professor,\\
        Ph.D.\ Ramis Akhmedov\\[8mm]
        \rule{6cm}{0.4pt}\\
        \textguillemotleft\rule{0.6cm}{0.4pt}\textguillemotright
        \rule{4cm}{0.4pt} 2026~y.
    \end{flushright}

    \vspace{20mm}
    \begin{center}
        {\Huge \textbf{DIPLOMA PROJECT}}\\[12mm]
        {\Large Oqyrman: A Web-Based Digital Library and\\
        Reservation Platform with an Integrated\\
        E-Reader and AI Reading Assistant}\\[10mm]
        {\large 6B06101 --- ``Information Systems''}
    \end{center}

    \vspace{15mm}
    \begin{center}
    \begin{tabular}{ll}
        \textbf{Supervisor:}     & Alikhan N.\\
        \textbf{Co-supervisor:}  & Baurzhan B.\\
        \textbf{Students:}       & Gainolla Azat\\
                                 & Arginbekov Alisher\\
                                 & Kengesbek Zanggar\\
                                 & Nassikhov Alikhan\\
    \end{tabular}
    \end{center}

    \vfill
    \begin{center}\textbf{Kaskelen, 2026}\end{center}
\end{titlepage}

% =============================================================================
%  ABSTRACT (English)
% =============================================================================
\chapter*{Abstract}
\addcontentsline{toc}{chapter}{Abstract}

This diploma project presents the design, implementation, and evaluation of
\textit{Oqyrman}, a comprehensive web-based digital library platform aimed at
modernising public and academic library services in the Republic of Kazakhstan.
The system unifies three traditionally fragmented experiences --- physical-copy
reservation, in-browser reading of digital books, and AI-assisted reading
support --- into a single multilingual web application that serves Kazakh,
Russian, and English readers from one account. The platform is built as a
decoupled client--server system: a React 18 single-page application written in
TypeScript on the frontend, and a Go (Gin) REST API backed by PostgreSQL on
the backend, deployed via Docker Compose behind an nginx reverse proxy and
released through a GitHub Actions continuous-delivery pipeline. Authentication
is handled by a stateless JWT scheme with single-flight refresh-token rotation
and Google OAuth integration. The integrated reader supports both EPUB and PDF
through \texttt{epubjs} and \texttt{react-pdf}, persists reading progress as a
percentage-based position together with a CFI or page locator so that the
reading session resumes byte-accurately across devices. A selection-driven
overlay routes highlighted passages to a server-proxied OpenAI-compatible
language model (GPT-4o-mini) capable of explaining, summarising, and
translating the selection between the three supported languages. Reservation
workflows are anchored by per-copy QR tokens that librarians scan from a
mobile device to confirm pickup and return, eliminating ambiguity in the
offline branch of the workflow without requiring dedicated hardware.
Functional walk-throughs and exploratory testing demonstrate that the system
reliably supports the entire reader life-cycle for a catalogue representative
of a medium-sized municipal library, and that the architecture remains
horizontally extensible to additional libraries, branches, and content types.

\textbf{Keywords:} digital library, e-reader, EPUB, PDF, JWT authentication,
React, Go, PostgreSQL, AI reading assistant, QR-based reservation, Kazakhstan.

% =============================================================================
%  АҢДАТПА (Kazakh)
% =============================================================================
\selectlanguage{kazakh}
\chapter*{Аңдатпа}
\addcontentsline{toc}{chapter}{Аңдатпа}

Бұл дипломдық жоба Қазақстан Республикасының кітапханалық қызметтерін
жаңғыртуға бағытталған \textit{Oqyrman} веб-платформасын жобалау, әзірлеу
және бағалау нәтижелерін ұсынады. Жүйе бұрын бір-бірінен бөлек жүзеге
асырылатын үш тәжірибені --- физикалық кітапты алдын ала брондау, цифрлық
кітапты браузер ішінде оқу және жасанды интеллект көмегімен оқуға қолдау
көрсету --- бір көптілді веб-қосымшаға біріктіреді және оқырманды қазақ,
орыс, ағылшын тілдерінде бір ғана аккаунттан қызмет алуға мүмкіндік береді.
Платформа клиент-серверлік архитектура негізінде құрылған: алдыңғы жағы
TypeScript-те жазылған React 18 SPA, артқы жағы Go (Gin) REST API және
PostgreSQL дерекқоры. Жүйе Docker Compose және nginx арқылы өндіріске
шығарылады, ал жеткізу GitHub Actions автоматтандырылған құбыры арқылы
жүзеге асырылады. Аутентификация JWT таңбалауышының күй сақтамайтын схемасы,
жалғыз ағынды (single-flight) жаңарту стратегиясы және Google OAuth
интеграциясымен қамтамасыз етілген. Кірістірілген оқу құралы EPUB және PDF
форматтарын \texttt{epubjs} және \texttt{react-pdf} арқылы қолдайды, оқу
ілгерілеуі пайыздық позиция мен CFI немесе бет нөмірі түрінде сақталады.
Белгіленген үзінді сервер арқылы OpenAI-үйлесімді GPT-4o-mini моделіне
жіберіліп, түсіндіру, қысқаша мазмұндау және үш тілдің кез келген жұбы
арасында аударма жасау қызметтерін ұсынады. Брондау процесі әр данаға
тағайындалған QR таңбалауышымен жабылады: кітапханашы оны смартфон
арқылы сканерлеп алу мен қайтаруды растайды.

\textbf{Кілт сөздер:} цифрлық кітапхана, электронды оқу құралы, EPUB, PDF,
JWT аутентификациясы, React, Go, PostgreSQL, AI оқу көмекшісі.

% =============================================================================
%  АННОТАЦИЯ (Russian)
% =============================================================================
\selectlanguage{russian}
\chapter*{Аннотация}
\addcontentsline{toc}{chapter}{Аннотация}

В настоящей дипломной работе представлены проектирование, разработка и
оценка веб-платформы \textit{Oqyrman}, направленной на модернизацию
библиотечного обслуживания в Республике Казахстан. Система объединяет три
традиционно разрозненных сценария --- бронирование физического экземпляра,
чтение цифровых книг непосредственно в браузере и поддержку процесса чтения
средствами искусственного интеллекта --- в единое многоязычное
веб-приложение, обслуживающее казахско-, русско- и англоязычных читателей
из одного аккаунта. Платформа построена по клиент-серверной архитектуре:
одностраничное приложение на React 18 с TypeScript на стороне клиента и
REST API на Go (фреймворк Gin) с PostgreSQL на стороне сервера. Развёртывание
выполняется через Docker Compose за обратным прокси nginx, доставка
автоматизирована при помощи GitHub Actions. Аутентификация реализована на
основе JWT с single-flight ротацией refresh-токена и интеграцией Google
OAuth. Встроенный ридер поддерживает форматы EPUB и PDF через библиотеки
\texttt{epubjs} и \texttt{react-pdf}, сохраняет прогресс чтения в виде
процентной позиции и CFI или номера страницы, обеспечивая бит-точное
возобновление сессии на любом устройстве. Слой выделения текста
маршрутизирует выбранный фрагмент на серверный прокси к OpenAI-совместимой
языковой модели GPT-4o-mini, которая выполняет пояснение, реферирование и
перевод между тремя поддерживаемыми языками. Процесс бронирования
закрывается одноразовым QR-токеном, который библиотекарь сканирует со
смартфона для подтверждения выдачи и возврата, без необходимости
специализированного оборудования.

\textbf{Ключевые слова:} цифровая библиотека, электронная читалка, EPUB, PDF,
JWT, React, Go, PostgreSQL, AI-ассистент чтения, бронирование по QR.

\selectlanguage{english}

% =============================================================================
%  TABLE OF CONTENTS
% =============================================================================
\tableofcontents
\listoffigures
\addcontentsline{toc}{chapter}{List of Figures}
\listoftables
\addcontentsline{toc}{chapter}{List of Tables}

\clearpage
\pagenumbering{arabic}

% =============================================================================
%  CHAPTER 1 — INTRODUCTION
% =============================================================================
\chapter{Introduction}

\section{Problem Background}

Public, university, and municipal libraries in Kazakhstan operate under two
simultaneous, structurally opposed pressures. On the one hand, their physical
collections remain a primary source of academic and recreational reading
material, particularly in regional branches where a substantial portion of
the catalogue is in Kazakh and Russian and is not available through
international commercial platforms such as Amazon Kindle or Apple Books. On
the other hand, the day-to-day workflow of a contemporary reader has shifted
decisively toward mobile-first, on-demand digital experiences. The reader is
no longer prepared to walk into a branch in order to discover whether a copy
is available, nor to keep a paper bookmark in order to remember a chapter
number; the modern reader expects to search the catalogue from home, hold a
copy until a convenient pickup time, track personal progress across devices,
and consume long-form text in print or on screen with parity of features.

Most existing library information systems in the country address only a
fragment of this expectation. A typical regional deployment exposes the
catalogue through a legacy OPAC interface that mirrors the internal
bibliographic record rather than the reader's mental model, requires physical
presence to confirm reservations, and treats digital reading as an entirely
separate product, often outsourced to a third-party vendor whose user
account is not linked to the reader's library identity. The result is a
fragmented experience in which the reader must reconcile two or three
independent accounts --- one for catalogue search, one for digital reading,
one for community discovery --- none of which retain a unified history of
what has been read or reserved.

The fragmentation is not merely a user-experience inconvenience. It produces
measurable operational consequences. Librarians must answer reservation
questions over the telephone because there is no self-service channel.
Readers abandon physical reservations because the confirmation step is
opaque. Acquisitions decisions are made without the consolidated reading
history that a unified platform would expose. From the standpoint of digital
inclusion, the absence of a Kazakh-first reader application means that a
reader who prefers Kazakh-language interfaces is implicitly nudged toward
either Russian-language or English-language commercial alternatives, which
in turn weakens the demand signal for Kazakh-language digital titles.

\medskip
\noindent\textit{Target audience.} The platform is positioned for three
overlapping reader segments inside Kazakhstan, with secondary applicability
to the diaspora abroad:

\begin{itemize}[leftmargin=*]
    \item \textbf{Schoolchildren and university students (16--24 years old)}
    who require both physical textbooks and digital long-form reading,
    typically in two or three languages, and who are accustomed to
    mobile-first, app-grade user experiences in adjacent product categories.
    \item \textbf{Adult lifelong readers (25--55 years old)} who maintain a
    long-term relationship with a specific branch but increasingly prefer
    digital reading during commute or travel, and who require the
    continuity of a single reading history across formats.
    \item \textbf{Librarians and library staff} who require an operational
    surface for confirming pickups and returns, monitoring inventory at
    the branch, and observing reader-side workflow without leaving the
    front desk or installing additional desktop software.
\end{itemize}

\noindent\textit{Motivation.} The author team observed, over the course of
preliminary domain study, that each of the four desirable capabilities
(catalogue, e-reader, reservation, AI-assisted reading) is individually
solved by some product in the international or domestic market, but none
of the surveyed products presents the four as facets of one reader account.
Closing this integration gap with a domestic, Kazakh-first platform
became the principal driver of the project.

\section{Aims and Objectives}

The overall aim of the work is to design, implement, and evaluate a
production-ready web platform that unifies catalogue browsing, physical
reservation, in-browser reading, and AI-assisted reading support under a
single user account, with first-class support for the Kazakh, Russian, and
English languages, and with operational characteristics suitable for a
medium-sized municipal or university library.

\noindent
To achieve this aim, the following specific objectives were defined:

\begin{enumerate}[label=\arabic*., leftmargin=*]
    \item Conduct a comparative study of existing digital library and
    reading platforms, both international and domestic, identify the
    functional gap that justifies a new system, and codify that gap as a
    set of requirements.
    \item Design a normalised relational data model that captures
    bibliographic entities, library inventory, reservations, reading
    sessions, notes, reviews, AI conversations, and notifications under
    one schema, and materialise the schema through exactly two canonical
    migration files (\texttt{init} and \texttt{seed}).
    \item Implement a stateless REST API in Go that enforces JWT-based
    authentication with single-flight refresh-token rotation, role-based
    access for librarians, parameterised SQL access, and resource-level
    authorisation derived from the request context.
    \item Implement a single-page React application that consumes the API,
    presents a responsive interface in three languages, embeds an EPUB
    and PDF reader with persistent reading progress and notes, and
    surfaces an AI-assisted selection layer inside the reader.
    \item Integrate a server-proxied large language model accessible
    through the reader to enable on-the-fly explanation, summarisation,
    and translation of selected text fragments, with the language of
    explanation derived from the reader's interface preference.
    \item Implement a QR-based reservation confirmation workflow suitable
    for use by library staff with a standard mobile device, with no
    additional hardware investment.
    \item Containerise the entire system and deploy it through a
    continuous-delivery pipeline that rebuilds and restarts the
    production environment on each push to the main branch.
    \item Treat security as a cross-cutting concern aligned with the OWASP
    Top Ten, including bcrypt password hashing, short-lived access
    tokens, server-side refresh-token allowlist, parameterised queries,
    and TLS-everywhere transport.
\end{enumerate}

\section{Significance of the Work}

The platform addresses a practical, locally relevant problem with a system
whose architecture is intentionally portable. By delivering a unified
reader account, integrated reservation, and AI-assisted reading in a
single deployment, the project lowers the operational cost of digitising a
regional library and reduces the friction that an end user encounters when
transitioning between physical and digital materials.

The relational schema, REST contract, and reader components are designed to
be reused or extended; for example, the addition of new content types
(audiobooks, periodicals, scholarly articles) requires schema extension and
reader adapters but does not require revision of the authentication or
reservation subsystems. Because the AI subsystem speaks an OpenAI-compatible
chat-completion contract, the model can be swapped to an alternative
provider without code changes on the client. Because authentication is
stateless and the deployment is containerised, the system can be scaled
horizontally by running additional API replicas behind the existing nginx
front-end. From an academic standpoint, the work demonstrates an end-to-end
engineering process: requirements derived from a domain study, a pragmatic
technology selection grounded in production-proven components, a layered
architecture, a documented data model, and a reproducible deployment
pipeline.

% =============================================================================
%  CHAPTER 2 — LITERATURE REVIEW
% =============================================================================
\chapter{Literature Review}

\section{Overview of the Domain}

The domain of digital library platforms has matured through three broad
generations. The first generation, exemplified by classical OPAC systems
such as Koha and Evergreen, focused on cataloguing and circulation control,
with end-user interfaces that mirrored the internal record structure rather
than the reader's mental model. The second generation, driven by the mass
adoption of e-book formats and dedicated reading devices, foregrounded the
reading experience itself: platforms such as Amazon Kindle, Apple Books,
and the Adobe Digital Editions ecosystem optimised for distribution and
DRM-protected consumption of EPUB and proprietary formats. The third and
current generation re-aggregates the two preceding strands and adds
personalisation, social reading, and AI-assisted features. Representative
platforms include Goodreads (community shelving and reviews), the
OverDrive/Libby ecosystem (library-mediated lending of digital titles),
and recent generative-AI-enabled reading clients that summarise or explain
text on demand.

The 2003 UNESCO \textit{Charter on the Preservation of the Digital
Heritage} \cite{unesco2003} emphasises that digitisation of culturally
significant material must serve democratic access rather than enclosure;
the Europeana initiative \cite{europeana2020} and the Mukurtu CMS project
\cite{christen2012} are widely cited as concrete instances of that
principle. In the present work the same principle guides the design of the
catalogue surface: every reader is granted equal access to the digital
collection, with no per-title licensing barrier, and every bibliographic
record is presented in the reader's language of choice where translation
data is available.

\section{Analysis of Existing Solutions}

\subsection{International Platforms}

\textbf{OverDrive / Libby.} The most widely deployed model for
library-mediated digital lending: a library purchases a fixed number of
licences, and patrons borrow them through a companion application. The
model solves the contractual side of e-book lending but is heavily
decoupled from the catalogue of physical holdings; users are expected to
switch applications when transitioning between print and digital, and
Kazakh-language content is essentially absent from the catalogue.

\textbf{Goodreads.} Dominant for shelving, reviews, and community
discovery, but purely a metadata layer: it does not interact with either
physical or digital lending. A reader who maintains a Goodreads shelf must
nonetheless reconstitute it independently inside any library application
they use.

\textbf{Amazon Kindle / Apple Books.} Vendor-locked digital reading
ecosystems. They offer the highest-fidelity reader implementations on the
market and have begun to integrate AI features (Kindle's word-lookup and
translation, Apple's on-device summarisation). However, they neither
accept third-party catalogue data nor expose physical-copy reservation,
and Kazakh-language coverage is marginal.

\textbf{Library-of-Congress-style federated catalogues} (WorldCat,
HathiTrust). Powerful for academic discovery but oriented toward research
librarianship; they are not consumer reading applications.

\subsection{Domestic and Regional Platforms}

\textbf{Kazakh National Electronic Library (kazneb.kz).} Offers a
substantial corpus of digitised cultural-heritage materials but exposes
them through a reader interface that is essentially a server-rendered PDF
viewer. There is no concept of a personal account, reading progress, or
shelving, and the reading experience does not adapt to font size or
viewport.

\textbf{University OPAC deployments based on Koha and ABIS.} Provide
catalogue search and limited reservation workflows, but require library
staff to mediate every transaction manually, and the reader-facing surface
is dated.

\textbf{Commercial e-book aggregators (Litres, Bookmate).} Offer a strong
reader experience for the Russian-language market but are paid
subscription products without integration with public-library inventories,
and Kazakh-language coverage is small.

\textbf{Duolingo / language-learning platforms.} Not direct competitors,
but the literature on these platforms \cite{vesselinov2012,licorish2018}
informs the design of the AI-assisted reading layer in Oqyrman: short,
selection-driven prompts consistently outperform long-form lectures for
attention retention in mobile-first audiences.

\section{Competitive Analysis}

Table~\ref{tab:competitive} summarises a feature-level comparison of
representative platforms against the four functional pillars that motivate
the present work: a unified reader-centric account, in-browser reading of
standard formats, reservation of physical copies tied to specific library
branches, and AI-assisted reading.

\begin{table}[H]
\centering
\caption{Competitive analysis of existing platforms.}
\label{tab:competitive}
\small
\begin{tabularx}{\textwidth}{lXXXXX}
\toprule
\textbf{Platform} & \textbf{Unified account} & \textbf{In-browser reader} &
\textbf{Physical reservation} & \textbf{AI assistant} & \textbf{Kazakh
content} \\
\midrule
OverDrive / Libby & Partial & Yes & No & No & No \\
Goodreads        & Yes & No & No & No & Marginal \\
Kazneb.kz        & No & Limited (PDF) & No & No & Yes \\
Koha (OPAC)      & Yes (basic) & No & Yes & No & Mixed \\
Kindle           & Vendor-locked & Yes (web reader) & No & Partial & No \\
Litres / Bookmate & Yes (paid) & Yes & No & No & Marginal \\
\textbf{Oqyrman (this work)} & \textbf{Yes} & \textbf{Yes (EPUB + PDF)} &
\textbf{Yes} & \textbf{Yes} & \textbf{Yes} \\
\bottomrule
\end{tabularx}
\end{table}

\textit{[FIGURE PLACEHOLDER 2.1: Screenshot collage of the home pages of
OverDrive, Goodreads, kazneb.kz, Litres, and Koha for visual comparison.
Insert five browser screenshots arranged in a 2$\times$3 grid with the last
cell left empty or used for a legend.]}

\section{Limitations of Existing Systems}

Three structural limitations recur across the surveyed systems and
constitute the negative space that Oqyrman is positioned to fill.

\textit{Account fragmentation.} The reader rarely owns a single identity
that spans both the physical and digital sides of the same library, and
reading history is therefore irreparably split. Even within the digital
side, vendor-locked ecosystems prevent the reader from carrying their
shelf to another platform.

\textit{Language coverage.} The dominant commercial platforms have
minimal Kazakh-language content, while the domestic platforms have minimal
modern reader features. Consequently, a Kazakh-first reader is forced to
choose between linguistic fidelity and product polish.

\textit{Operational friction in reservation.} Where physical reservation
exists, confirmation is typically performed by manual entry into a
back-office client; there is no lightweight mechanism (such as a
scannable token) that allows a librarian to confirm pickup or return
without dedicated hardware.

\section{Gap Analysis}

Synthesising the above, the gap that the present work addresses is the
\textit{integration} gap rather than any single missing feature. Each
individual capability --- catalogue, e-reader, reservation, AI assistant
--- is available in isolation in some existing system, but none of the
surveyed solutions presents them as facets of one reader account, with
shared authentication, shared reading history, and shared bibliographic
identifiers, in a Kazakh-first interface, with a librarian-side workflow
deliverable on a standard smartphone. Oqyrman is positioned to fill
exactly this integration gap, with the additional design constraint that
every reader-facing surface (catalogue, reader, AI overlay, notifications)
must be available in Kazakh, Russian, and English without a reload.

% =============================================================================
%  CHAPTER 3 — DESIGN AND METHODOLOGY
% =============================================================================
\chapter{Design and Methodology}

\section{Methodology}

\subsection{Development Process}

The project was developed under an iterative, Scrum-inspired methodology
adapted to the constraints of a four-author bachelor-level undertaking.
Work was organised into one-week iterations (sprints). Each sprint opened
with a planning meeting that enumerated the user stories targeted for the
iteration and closed with a retrospective and demo against the acceptance
criteria of those stories. The product backlog was maintained as a flat
list of user stories on Jira, each tagged with one of three priority
bands: \textit{must} (required for the minimum viable release),
\textit{should} (required for a complete experience), and \textit{could}
(deferred enhancements). All source code was tracked in a Git repository
hosted on GitHub, with feature work performed on short-lived branches and
merged into \texttt{main} only after the corresponding peer code review.

\textit{[FIGURE PLACEHOLDER 3.1: Scrum process diagram showing the
sequence Vision $\to$ Product Backlog $\to$ Sprint Planning $\to$ Sprint
Backlog $\to$ 1-week Sprint $\to$ Daily Standup $\to$ Sprint Review $\to$
Retrospective. Insert as a horizontal infographic.]}

\textit{[FIGURE PLACEHOLDER 3.2: Screenshot of the Jira board for the
Oqyrmandar team showing TODO / IN-PROGRESS / DONE columns with
representative tickets such as ``Implement JWT refresh middleware'',
``Build EPUB reader settings panel'', ``Stream AI responses via SSE''.]}

\subsection{Team Roles}

The four-person team distributed responsibilities along the natural
boundaries of the system, with cross-coverage on overlapping areas.
Table~\ref{tab:roles} summarises the assignment.

\begin{table}[H]
\centering
\caption{Team roles within the Oqyrmandar project.}
\label{tab:roles}
\small
\begin{tabularx}{\textwidth}{l X}
\toprule
\textbf{Member} & \textbf{Primary responsibility} \\
\midrule
Nassikhov Alikhan  & Project Manager / Team Lead. Backlog grooming,
sprint planning, stakeholder communication, cross-component integration. \\
Arginbekov Alisher & Backend Developer. Go/Gin REST API, JWT
authentication, PostgreSQL schema, SQL repositories, AI server-side
adapter. \\
Kengesbek Zanggar  & Frontend Developer. React 18 SPA, shadcn/ui
components, EPUB and PDF reader, AI selection layer, i18n surface. \\
Gainolla Azat      & DevOps Engineer. Docker images, Docker Compose
production stack, nginx configuration, GitHub Actions CI/CD, deployment
to the production VM. \\
\bottomrule
\end{tabularx}
\end{table}

\subsection{Quality Practices}

Several practices were institutionalised across the team to maintain
quality at the speed of weekly delivery. \textit{Code review} was
mandatory before any merge to \texttt{main}: every pull request required
at least one approving review, and the reviewer was expected to
specifically check security-sensitive changes (authentication, SQL,
external HTTP calls). \textit{Automated tests} were executed by GitHub
Actions on every push: the frontend used Vitest and
\texttt{@testing-library/react}, the backend used Go's standard
\texttt{testing} package. \textit{Continuous deployment}: a successful
merge to \texttt{main} automatically rebuilt the Docker images and
restarted the production stack via SSH, ensuring that the end-to-end
pipeline was exercised at the cadence of the development cycle rather
than only at release boundaries. \textit{Heuristic UX evaluation}
following Nielsen's classic ten heuristics \cite{nielsen1994} was
performed by the team at each sprint review, with discovered issues
re-entered into the backlog.

\section{System Design}

\subsection{Visual Identity and Colour Palette}

A consistent visual identity was developed before any production-grade
component was built, in order to anchor subsequent design decisions to a
shared aesthetic language. The palette draws from warm, paper-adjacent
neutrals appropriate to a reading product, accented by a single saturated
primary that echoes the warm tones of traditional Kazakh ornamentation
without being literal. The palette is defined in CSS custom properties in
\texttt{src/index.css} and consumed throughout the component tree via
Tailwind's semantic tokens, which makes dark-mode and theming a single
source of truth.

\textit{[FIGURE PLACEHOLDER 3.3: Colour palette swatches: primary
(saturated warm orange), surface (warm cream), foreground (near-black),
muted (warm grey), accent (deep teal). Display as five circles or
rectangles labelled with HEX values.]}

\subsection{System Architecture}

The system follows a classical three-tier decoupled architecture: a
stateless single-page-application client, a stateless REST API server,
and a relational database. A separate object-storage volume hosts media
assets (book covers, EPUB and PDF files), exposed to the client through
public URLs served by the same nginx instance that fronts the SPA.

\textit{[FIGURE PLACEHOLDER 3.4: High-level architecture diagram showing
the following components and arrows. Components: Browser (React SPA),
nginx (reverse proxy plus static files), Go/Gin API container,
PostgreSQL container, file-storage volume. External services: Google
OAuth, OpenAI-compatible LLM endpoint (GPT-4o-mini), SMTP/email provider.
Arrows: Browser $\to$ nginx; nginx $\to$ Go API at \texttt{/api/v1/*};
nginx $\to$ static SPA; Go API $\to$ PostgreSQL; Go API $\to$ Email; Go
API $\to$ LLM (server-proxied); Browser $\to$ Google OAuth (for
\texttt{id\_token} only).]}

The frontend is delivered as a fully static bundle produced by Vite (with
SWC) and served from nginx with SPA fallback rewrites; it never executes
any server-side rendering. All authenticated state is materialised on the
client from a JWT held in \texttt{localStorage}, augmented by a
per-request \texttt{Authorization: Bearer} header. The backend is a
single Go binary compiled into a Docker image; it owns all business
logic, validates tokens, and is the only component allowed to read or
write the database. The LLM-backed AI assistant is implemented as a
server-side adapter: the client never holds credentials for the
third-party model, and all prompt assembly, redaction of personally
identifiable information, and rate-limiting are performed in the API tier.

\subsection{Database Design}

The database is a normalised PostgreSQL schema. The entity-relationship
model is centred on the \textit{book} entity, which acts as the canonical
bibliographic record. Around it the schema layers ownership (library
inventories), social signal (reviews, wishlists), behavioural state
(reading sessions, notes), workflow (reservations, notifications), and
conversational context (AI conversations and messages).

\begin{itemize}[leftmargin=*]
    \item \textbf{users} --- canonical identity record (\texttt{id, name,
    surname, email, phone, avatar\_url, role, password\_hash,
    has\_password}). The \texttt{role} field distinguishes ordinary
    readers from librarians.
    \item \textbf{authors} and \textbf{genres} --- reference dictionaries
    for bibliographic metadata; \texttt{authors} carries bilingual
    biographies.
    \item \textbf{books} --- the bibliographic record, joined to one
    \texttt{author} and one \texttt{genre}, carrying ISBN, year,
    \texttt{total\_pages}, language, and aggregate rating columns
    (\texttt{avg\_rating}, \texttt{ratings\_count}).
    \item \textbf{book\_files} --- 1:N relationship attaching one or more
    digital files (EPUB, PDF) to a book; the reader chooses an adapter
    based on the \texttt{format} column.
    \item \textbf{libraries} --- physical branches with geographic
    coordinates consumed by the Leaflet map view.
    \item \textbf{library\_books} --- join table carrying
    \texttt{total\_copies} and \texttt{available\_copies} for each
    (\texttt{library}, \texttt{book}) pair; reservations decrement and
    restore these counters atomically.
    \item \textbf{reservations} --- workflow record with \texttt{status}
    in \{\texttt{pending, active, completed, cancelled}\}, due-date
    tracking, extension counter, and a one-time \texttt{qr\_token} used
    for librarian confirmation.
    \item \textbf{wishlist\_items} --- reader's personal shelves with the
    \texttt{status} discriminator \{\texttt{want\_to\_read, reading,
    finished}\}.
    \item \textbf{reading\_sessions} --- per-(\texttt{user, book})
    progress record holding a \texttt{progress} percentage and a
    \texttt{cfi\_position} (EPUB CFI string or PDF page locator) so that
    the reader can resume across devices.
    \item \textbf{reading\_notes} --- highlight-anchored textual notes;
    each note carries a serialised position descriptor that the reader
    uses to re-anchor the highlight.
    \item \textbf{reviews} --- public review and rating per (\texttt{user,
    book}); a database trigger maintains the aggregate rating columns on
    the parent \texttt{books} row.
    \item \textbf{notifications} --- in-app notifications with a typed
    \texttt{type} discriminator (e.g.\ \texttt{reservation\_success},
    \texttt{return\_deadline}, \texttt{event\_reminder}).
    \item \textbf{events} --- bilingual library events surfaced on a
    dedicated page.
    \item \textbf{ai\_conversations} and \textbf{ai\_messages} ---
    threaded storage of AI dialogues, with \texttt{role} discriminator
    \{\texttt{user, assistant}\} on each message.
\end{itemize}

\textit{[FIGURE PLACEHOLDER 3.5: Entity-relationship diagram of the
schema above. Render with dbdiagram.io or pgAdmin's ERD view; place
\texttt{users} on the left, \texttt{books} in the centre, and the
workflow tables (\texttt{reservations}, \texttt{reading\_sessions},
\texttt{wishlist\_items}, \texttt{ai\_conversations}) radiating outward.
Export as PNG, full-width.]}

In line with the project's migration convention, the schema is
materialised through exactly two migration files: an \texttt{init}
migration that defines all tables, indexes, and triggers, and a
\texttt{seed} migration that loads demonstration data. Schema or seed
changes are folded into these two files rather than appended as numbered
increments; this keeps the canonical schema unambiguous and avoids long
migration chains. Listing~\ref{lst:init} shows a representative excerpt
from the \texttt{init} migration.

\begin{lstlisting}[language=SQL,caption={Excerpt of the \texttt{init} migration: \texttt{books} and \texttt{reservations}.},label=lst:init]
CREATE TABLE books (
    id              BIGSERIAL PRIMARY KEY,
    title           TEXT NOT NULL,
    description     TEXT,
    description_kk  TEXT,
    isbn            TEXT,
    year            INT,
    total_pages     INT,
    language        TEXT NOT NULL,
    cover_url       TEXT,
    author_id       BIGINT REFERENCES authors(id) ON DELETE RESTRICT,
    genre_id        BIGINT REFERENCES genres(id)  ON DELETE RESTRICT,
    avg_rating      NUMERIC(3,2) NOT NULL DEFAULT 0,
    ratings_count   INT          NOT NULL DEFAULT 0,
    created_at      TIMESTAMPTZ  NOT NULL DEFAULT now()
);

CREATE TABLE reservations (
    id                BIGSERIAL PRIMARY KEY,
    user_id           BIGINT NOT NULL REFERENCES users(id),
    library_book_id   BIGINT NOT NULL REFERENCES library_books(id),
    status            TEXT   NOT NULL CHECK (status IN
                          ('pending','active','completed','cancelled')),
    reserved_at       TIMESTAMPTZ NOT NULL DEFAULT now(),
    due_date          DATE,
    returned_at       TIMESTAMPTZ,
    extended_count    INT  NOT NULL DEFAULT 0,
    qr_token          TEXT UNIQUE
);
\end{lstlisting}

\subsection{Data Flow}

A representative interaction --- the reservation of a physical book ---
is illustrative of the data-flow pattern that the system follows
throughout.

\begin{enumerate}[leftmargin=*]
    \item The client issues \texttt{GET /books/\{id\}} to fetch the
    bibliographic record together with embedded \texttt{author},
    \texttt{genre}, and \texttt{file} relations.
    \item The client issues \texttt{GET /books/\{id\}/availability} to
    fetch the per-library inventory rows; the API joins
    \texttt{library\_books} with \texttt{libraries} and returns the
    available copy count for each branch.
    \item The reader selects a branch and submits
    \texttt{POST /reservations} with the chosen
    \texttt{library\_book\_id}. The server enforces, in a single
    transaction, that \texttt{available\_copies > 0}, decrements the
    counter, inserts a \texttt{reservations} row in status
    \texttt{pending}, and emits a \texttt{reservation\_success}
    notification.
    \item The server returns the reservation, including a freshly minted
    one-time \texttt{qr\_token}; the client renders the token as a QR
    code via \texttt{qrcode.react}.
    \item At the library, the librarian's account scans the QR code
    through the in-app HTML5 scanner (\texttt{html5-qrcode}); the scanner
    client resolves the token via \texttt{POST /reservations/scan},
    which advances the reservation to status \texttt{active} and emits a
    \texttt{pickup\_success} notification.
    \item Return is handled symmetrically; the server marks the
    reservation \texttt{completed}, restores the
    \texttt{available\_copies} counter, and emits a final notification.
\end{enumerate}

\textit{[FIGURE PLACEHOLDER 3.6: BPMN-style data-flow diagram showing
three swimlanes (Reader, API/WMS, Librarian) and the messages exchanged
between them across the reservation life-cycle.]}

The same general pattern --- client read, optimistic UI rendering, atomic
write through the API, server-side notification emission --- governs the
wishlist, review, and reading-session subsystems.

\section{Implementation Details}

\subsection{Backend (Go / Gin)}

The backend is implemented in Go using the Gin HTTP framework. The
codebase is organised by resource (\texttt{auth}, \texttt{books},
\texttt{libraries}, \texttt{reservations}, \texttt{ai}, $\dots$), with
each resource exposing a handler module, a service module, and a
repository module. The repository layer wraps a connection pool to
PostgreSQL; SQL is written by hand rather than generated by an ORM, which
keeps query plans predictable and lets the schema express constraints
(foreign keys, partial indexes) directly. Authentication is handled by a
dedicated middleware that validates the incoming JWT, resolves the
associated \texttt{users} row, and injects it into the request context;
downstream handlers read the authenticated user from the context and
never trust client-supplied identifiers.

Listing~\ref{lst:jwt} shows the core of the JWT middleware. Refresh-token
rotation is implemented with a server-side allowlist: each issued refresh
token is recorded in the database, single-use, and invalidated upon
redemption, preventing replay.

\begin{lstlisting}[language=Go,caption={JWT authentication middleware (excerpt).},label=lst:jwt]
func RequireAuth(secret []byte, repo UserRepo) gin.HandlerFunc {
    return func(c *gin.Context) {
        h := c.GetHeader("Authorization")
        if !strings.HasPrefix(h, "Bearer ") {
            c.AbortWithStatusJSON(401, gin.H{"code": "unauthorized"})
            return
        }
        tok, err := jwt.Parse(strings.TrimPrefix(h, "Bearer "),
            func(t *jwt.Token) (interface{}, error) { return secret, nil })
        if err != nil || !tok.Valid {
            c.AbortWithStatusJSON(401, gin.H{"code": "invalid_token"})
            return
        }
        claims := tok.Claims.(jwt.MapClaims)
        uid := int64(claims["sub"].(float64))
        user, err := repo.ByID(c.Request.Context(), uid)
        if err != nil {
            c.AbortWithStatusJSON(401, gin.H{"code": "user_not_found"})
            return
        }
        c.Set("user", user)
        c.Next()
    }
}
\end{lstlisting}

\subsection{Frontend (React 18 / Vite)}

The frontend is a React 18 single-page application written in TypeScript
and bundled with Vite (SWC). The component library combines Radix UI
primitives (through \texttt{shadcn/ui}) with Tailwind CSS for utility
styling. Server state is managed by TanStack Query: queries are keyed by
resource and filter parameters, mutations invalidate the affected keys,
and a global query client is configured with conservative stale times to
avoid redundant refetches during navigation. Forms are implemented with
\texttt{react-hook-form} and validated by \texttt{zod} schemas, which
doubles as type derivation for the form payload.

Routing is handled by React Router v6 in three tiers, each wrapped only
in the chrome it requires: a \textit{public} tier for the landing page
and authentication flows, a \textit{full-screen authenticated} tier for
the reader (which deliberately omits the navbar and footer to maximise
reading area), and an \textit{authenticated app} tier for catalogue,
profile, and workflow pages. Lazy loading via \texttt{React.lazy} and
\texttt{Suspense} defers each tier's bundle until first navigation.

The single-source-of-truth API client is centralised in
\texttt{src/lib/api/client.ts}; it injects the access token, intercepts
401 responses, and serialises concurrent refreshes into a single
in-flight refresh promise. Listing~\ref{lst:apiclient} reproduces the
single-flight refresh portion.

\begin{lstlisting}[language=JavaScript,caption={Single-flight refresh in the frontend API client (excerpt).},label=lst:apiclient]
let isRefreshing = false;
let refreshQueue: Array<(token: string | null) => void> = [];

async function refreshOnce(): Promise<string | null> {
  if (isRefreshing) {
    return new Promise(resolve => refreshQueue.push(resolve));
  }
  isRefreshing = true;
  try {
    const r = await fetch('/api/v1/auth/refresh', {
      method: 'POST',
      body: JSON.stringify({ refresh_token: tokenStorage.refresh() }),
    });
    if (!r.ok) throw new Error('refresh_failed');
    const { access_token, refresh_token } = await r.json();
    tokenStorage.set(access_token, refresh_token);
    refreshQueue.forEach(cb => cb(access_token));
    return access_token;
  } catch {
    notifySessionExpired();
    refreshQueue.forEach(cb => cb(null));
    return null;
  } finally {
    refreshQueue = [];
    isRefreshing = false;
  }
}
\end{lstlisting}

\subsection{The Integrated Reader}

The reader is composed of two format-specific components ---
\texttt{EpubReader} backed by \texttt{epubjs} and \texttt{PdfReader}
backed by \texttt{react-pdf} --- behind a common shell that owns settings
(font, font size, theme, line height), notes, progress synchronisation,
and the AI-selection layer. Reading progress is tracked as a percentage
rather than a page number; this is necessary because EPUB content reflows
according to the viewport and font settings, so a fixed page mapping
would be unstable across sessions. The \texttt{cfi\_position} string (or
page index for PDF) is persisted alongside the percentage so that
resumption is byte-accurate.

A custom hook abstracts the persistence concern from the format adapters
(Listing~\ref{lst:reader}). The hook debounces writes to avoid spurious
network traffic during rapid scrolling and reconciles optimistic local
state with server-confirmed state.

\begin{lstlisting}[language=JavaScript,caption={\texttt{useReadingSession} hook (excerpt).},label=lst:reader]
export function useReadingSession(bookId: string) {
  const [progress, setProgress] = useState(0);
  const [cfi, setCfi] = useState<string | null>(null);

  useEffect(() => {
    readingSessionsApi.get(bookId).then(s => {
      setProgress(s.progress); setCfi(s.cfi_position ?? null);
    });
  }, [bookId]);

  const persist = useMemo(
    () => debounce((p: number, c: string | null) =>
      readingSessionsApi.upsert({ book_id: bookId, progress: p,
                                  cfi_position: c }),
      1500),
    [bookId]
  );

  const update = useCallback((p: number, c: string | null) => {
    setProgress(p); setCfi(c); persist(p, c);
  }, [persist]);

  return { progress, cfi, update };
}
\end{lstlisting}

\subsection{AI Reading Assistant}

The AI subsystem is implemented as a streaming endpoint that proxies to
an OpenAI-compatible chat-completion API (GPT-4o-mini in the production
deployment). The server assembles the prompt from (i) a system message
that constrains the assistant to reading-related tasks in the user's
chosen interface language, (ii) the last $N$ messages of the
conversation truncated to a token budget, and (iii) the optional
selection passed by the reader (a paragraph excerpt with a directive
such as ``explain'', ``summarise'', or ``translate''). Tokens are
streamed back to the client as Server-Sent Events, and both the user
message and the assistant message are persisted to the
\texttt{ai\_messages} table once the stream concludes.

The assistant is reachable from two surfaces. The first is a global
floating chat widget mounted at the application layout level, suitable
for catalogue-wide queries (``recommend a Kazakh-language novel set in
the steppe''). The second is a contextual selection layer inside the
reader: highlighting a paragraph reveals an inline popover with three
primary actions, and the response is streamed in a side panel that does
not displace the reading area.

\subsection{Authentication and Authorisation}

The authentication subsystem implements four flows: email-and-password
registration with a 6-digit verification code, password-based login,
Google sign-in via the \texttt{@react-oauth/google} provider, and
password reset via a 6-digit code rather than an email link. Access
tokens are short-lived (15 minutes); refresh tokens live in a
server-side allowlist with single-use semantics, eliminating the entire
class of replay attacks against long-lived refresh tokens. Authorisation
is performed by reading the authenticated user from the request context
populated by the JWT middleware; user-supplied identifiers are never
trusted as authorisation. Librarian-only endpoints (notably
\texttt{POST /reservations/scan}) check the \texttt{role} column.

\textit{[FIGURE PLACEHOLDER 3.7: Sequence diagram of the JWT refresh
flow showing the browser, the API, and the database. Highlight the
single-flight behaviour: two concurrent 401 responses from protected
endpoints, both queued behind one outstanding
\texttt{POST /auth/refresh}.]}

\subsection{API Surface}

The HTTP API is REST-shaped and versioned under the \texttt{/api/v1}
prefix. Resources are addressed in the customary plural form,
collection endpoints accept a uniform pagination contract
(\texttt{?limit=\&offset=}), and the response body for a paginated
resource is \texttt{\{ items, total, limit, offset \}}. Errors are
returned as JSON objects with a stable \texttt{code} field and a
human-readable \texttt{message}; the client surfaces these through a
typed \texttt{ApiException} class. Table~\ref{tab:endpoints} catalogues
a representative subset.

\begin{table}[H]
\centering
\caption{Selected API endpoints.}
\label{tab:endpoints}
\small
\begin{tabularx}{\textwidth}{l l X}
\toprule
\textbf{Method} & \textbf{Path} & \textbf{Purpose} \\
\midrule
POST   & /auth/register             & Create account; sends 6-digit verification code \\
POST   & /auth/verify-email         & Confirm code, return access + refresh tokens \\
POST   & /auth/login                & Email + password authentication \\
POST   & /auth/google               & Exchange Google \texttt{id\_token} for session \\
POST   & /auth/refresh              & Rotate refresh token, return new pair \\
POST   & /auth/forgot-password      & Send 6-digit reset code \\
POST   & /auth/reset-password       & Consume reset code, set new password \\
GET    & /books                     & Catalogue with filters \texttt{q, genre\_id, author\_id, language} \\
GET    & /books/\{id\}              & Single book + relations + file \\
GET    & /books/\{id\}/availability & Per-library inventory \\
GET    & /libraries                 & List branches (with coordinates) \\
POST   & /reservations              & Create reservation, mint QR token \\
POST   & /reservations/scan         & Librarian: confirm pickup or return by QR \\
GET    & /wishlist                  & Reader shelves \\
POST   & /reading-sessions          & Upsert reading progress \\
GET    & /notes                     & List notes for a book \\
POST   & /ai/conversations          & Start a new AI conversation \\
POST   & /ai/conversations/\{id\}/messages & Stream assistant response (SSE) \\
GET    & /notifications             & Reader inbox \\
\bottomrule
\end{tabularx}
\end{table}

\subsection{Deployment and CI/CD}

The production deployment runs on a single virtual machine behind nginx,
which terminates TLS and reverse-proxies the API container while serving
the SPA bundle from disk. The full stack is described by a
\texttt{docker-compose.prod.yml} file that defines four services
(frontend, backend, postgres, and a one-shot migrator), with named
volumes for the database and the media storage. A push to \texttt{main}
triggers a GitHub Actions workflow that SSHes to the server and runs
\texttt{docker compose -f docker-compose.prod.yml up -d --build}, which
rebuilds only the changed images and restarts the corresponding
containers. The frontend's \texttt{VITE\_API\_BASE\_URL} is baked in at
build time via a Docker \texttt{ARG}; changing it for production is a
rebuild rather than a config reload, which is intentional --- it keeps
the deployed bundle deterministic.

\section{Security Considerations}

Security was treated as a cross-cutting design constraint. The platform
addresses each of the relevant items of the OWASP Top Ten 2021
\cite{owasp2021} as follows.

\begin{itemize}[leftmargin=*]
    \item \textbf{Broken authentication.} Passwords are hashed with
    bcrypt at a tunable cost factor; access tokens are short-lived (15
    minutes); refresh tokens are stored server-side in an allowlist with
    single-use semantics. The frontend serialises concurrent
    401-driven refreshes into a single in-flight refresh promise,
    eliminating the race condition that would otherwise produce
    duplicate refresh requests.
    \item \textbf{Broken access control.} Every authenticated handler
    reads the caller's identity from the request context populated by
    the JWT middleware; user-supplied identifiers are never trusted.
    Librarian-only endpoints check the \texttt{role} column.
    \item \textbf{Injection.} All database access uses parameterised
    queries; user-controlled strings are never concatenated into SQL.
    \item \textbf{Cross-site scripting.} The React renderer escapes
    interpolated values by default; the reader does not render
    user-controlled HTML, only sanitised EPUB and PDF content rendered
    by vetted libraries (\texttt{epubjs}, \texttt{react-pdf}).
    \item \textbf{Sensitive data exposure.} TLS terminates at the nginx
    layer; HTTP traffic is redirected to HTTPS. Tokens are never logged.
    \item \textbf{Cross-site request forgery.} Because authentication is
    carried by an explicit \texttt{Authorization} header rather than by
    a cookie, the API is intrinsically CSRF-resistant for cross-origin
    clients; a strict CORS allowlist limits the set of permitted
    origins.
    \item \textbf{Insecure deserialisation.} Request bodies are bound
    through Gin's strict JSON binding, which rejects unknown fields
    where annotated.
    \item \textbf{Security logging and monitoring failures.} The Go
    binary emits structured JSON logs that are aggregated to the host
    journal; failed authentications, refresh-token misuse attempts, and
    librarian scan operations are logged with correlation identifiers
    sufficient for after-the-fact investigation.
\end{itemize}

\section{Tools and Environment}

The technology stack reflects the principle of preferring widely
supported, production-proven components over novel but unproven ones.
Table~\ref{tab:stack} catalogues the principal tools.

\begin{table}[H]
\centering
\caption{Tools and technologies used.}
\label{tab:stack}
\small
\begin{tabularx}{\textwidth}{l X}
\toprule
\textbf{Layer} & \textbf{Technology} \\
\midrule
Frontend framework      & React 18, TypeScript \\
Build tool              & Vite (with SWC) \\
UI primitives           & Radix UI via shadcn/ui, Tailwind CSS \\
Server state            & TanStack Query \\
Forms / validation      & react-hook-form, zod \\
Routing                 & React Router v6 \\
Reader engines          & epubjs (EPUB), react-pdf (PDF) \\
Mapping                 & Leaflet, react-leaflet \\
QR codes                & qrcode.react, html5-qrcode \\
Internationalisation    & i18next, react-i18next (en, ru, kk) \\
Backend language        & Go (Gin) \\
Database                & PostgreSQL \\
Authentication          & JWT (custom), Google OAuth \\
LLM integration         & OpenAI-compatible chat completion (GPT-4o-mini, server-proxied) \\
Containerisation        & Docker, Docker Compose \\
Reverse proxy / TLS     & nginx \\
CI/CD                   & GitHub Actions (SSH deploy) \\
Testing                 & Vitest, @testing-library/react, jsdom, Go testing \\
Project tracking        & Jira, GitHub Issues \\
Design tooling          & Figma, Miro \\
API testing             & Postman, Swagger UI \\
\bottomrule
\end{tabularx}
\end{table}

% =============================================================================
%  CHAPTER 4 — RESULTS AND DISCUSSION
% =============================================================================
\chapter{Results and Discussion}

\section{Implemented Features}

The delivered system covers the full reader-facing surface defined in the
objectives chapter. The following subsections walk through the principal
implemented features and the screens that surface them.

\subsection{Authentication and Onboarding}

Email-and-password registration is gated by a 6-digit verification code
delivered to the supplied address, after which the API issues an access
and a refresh token and the client persists them via \texttt{tokenStorage}.
Password-based login and Google sign-in share the same token-issuance
path. Password reset is performed via a 6-digit reset code rather than an
email link, which avoids a class of phishing concerns associated with
clickable reset links and produces a self-contained mobile-friendly flow.

\textit{[FIGURE PLACEHOLDER 4.1: Screenshot of the login screen showing
email and password fields, ``Sign in with Google'' button, and
``Forgot password?'' link.]}

\textit{[FIGURE PLACEHOLDER 4.2: Screenshot of the 6-digit code entry
component used by both \texttt{verify-email} and \texttt{reset-password}
flows.]}

\subsection{Catalogue and Discovery}

The catalogue component implements full-text search across title,
author, and ISBN, combined with multi-faceted filtering by genre,
author, language, and publication-year range; results are paginated
through the standard \texttt{limit/offset} contract and rendered as
responsive book cards with cover, rating, and authoring metadata. The
book detail page joins the bibliographic record with per-library
inventory, reviews, and a similar-books recommendation row computed
server-side from genre and author overlap.

\textit{[FIGURE PLACEHOLDER 4.3: Screenshot of the catalogue page with
active filters (e.g.\ language = Kazakh, genre = Modern fiction),
showing the responsive grid of book cards on a desktop viewport.]}

\textit{[FIGURE PLACEHOLDER 4.4: Screenshot of a book detail page (e.g.\
``Абай жолы'' by Мұхтар Әуезов) displaying the cover, metadata, action
buttons (Read / Reserve / Add to wishlist), the per-library availability
table, and the similar-books row.]}

\subsection{Physical Reservation}

The reservation workflow operates end-to-end. A reader selects a
library, submits the reservation, and is presented with a QR code that
they bring to the branch. The librarian, signed in under a
librarian-role account, opens the in-app scanner; on a successful scan
the system advances the reservation status and decrements or restores
the inventory atomically inside a single SQL transaction. The
reservation detail page surfaces the due date, the extension counter,
and the action buttons appropriate to the current status.

\textit{[FIGURE PLACEHOLDER 4.5: Screenshot of an active reservation
card showing the QR code, the due date, the ``Cancel'' / ``Extend''
actions, and the branch address.]}

\textit{[FIGURE PLACEHOLDER 4.6: Screenshot of the librarian-side
scanner view (HTML5 camera viewport) with a successful-scan toast.]}

\subsection{Reading Experience}

The integrated reader supports both EPUB and PDF inputs. Settings (font
family, font size, line height, theme) are persisted per user; reading
progress and the resumption pointer are persisted per (\texttt{user,
book}) pair through the reading-sessions endpoint, which is debounced
on the client at 1.5 s to avoid spurious writes during rapid scrolling.
The note layer allows the reader to highlight a passage and attach a
textual note; the note's anchor is serialised so that it survives reflow.

\textit{[FIGURE PLACEHOLDER 4.7: Screenshot of the EPUB reader in dark
mode displaying a Kazakh-language chapter, with a highlighted passage
and the note editor sidebar open on the right.]}

\textit{[FIGURE PLACEHOLDER 4.8: Screenshot of the reader settings
panel: font family, font size slider, line height, and theme toggle.]}

\subsection{AI Reading Assistant}

The AI assistant is reachable both as a global widget (a floating chat
launcher mounted at the application layout level) and as a contextual
overlay inside the reader, where the reader can select a paragraph and
request an explanation, summary, or translation. Conversations are
listed in a sidebar and can be resumed; messages stream in
token-by-token through the SSE endpoint, with a stop button that
cancels the upstream HTTP request immediately on press.

\textit{[FIGURE PLACEHOLDER 4.9: Screenshot of the reader's AI
selection overlay: a paragraph is highlighted and a small popover
offers ``Explain'', ``Summarise'', ``Translate''. The adjacent panel
shows the streamed assistant reply.]}

\textit{[FIGURE PLACEHOLDER 4.10: Screenshot of the global AI chat
widget on the catalogue page with a sample prompt (``recommend a
Kazakh novel set in the steppe'') and a streamed response.]}

\subsection{Wishlist, Reviews, and Notifications}

The wishlist supports three shelves --- \textit{Want to read},
\textit{Reading}, and \textit{Finished} --- with status-dropdown
transitions that immediately update the server-side state and
optimistically update the client-side cache. Reviews are public and
constrained to one review per (\texttt{user, book}); a database
trigger maintains the aggregate \texttt{avg\_rating} and
\texttt{ratings\_count} columns on the parent \texttt{books} row, so
the catalogue surface never has to recompute aggregates at read time.
The notifications inbox surfaces the typed notifications emitted by
the workflow subsystems with appropriate icons and deep-links into
the related resource.

\textit{[FIGURE PLACEHOLDER 4.11: Screenshot of the wishlist page
showing the three-shelf tab bar and the cards in the ``Reading''
shelf.]}

\textit{[FIGURE PLACEHOLDER 4.12: Screenshot of the notifications
inbox with a mix of \texttt{reservation\_success} and
\texttt{return\_deadline} entries.]}

\subsection{Libraries Map and Events}

The libraries page renders a Leaflet map of all branches together with
a list view; clicking a marker scrolls the corresponding card into
view. Events are published bilingually and surfaced on a dedicated
page with client-side filtering by upcoming-vs-past status.

\textit{[FIGURE PLACEHOLDER 4.13: Screenshot of the libraries page
with the Leaflet map and the per-branch cards.]}

\textit{[FIGURE PLACEHOLDER 4.14: Screenshot of the reader's
profile page showing avatar, statistics (books read, reading-time
hours), and the language switcher.]}

\section{System Capabilities}

The aggregate capabilities of the platform can be summarised along
five axes:

\begin{itemize}[leftmargin=*]
    \item \textbf{Multilingual presentation.} Every reader-facing
    string is routed through \texttt{i18next} with English, Russian,
    and Kazakh translation tables; the language switcher persists the
    selection and triggers a no-reload re-render. Bibliographic
    free-text fields (description, biography, event description)
    carry both Russian and Kazakh variants in the database, allowing
    language-aware delivery at the data layer rather than purely at
    the UI layer.
    \item \textbf{Cross-device session continuity.} The combination of
    server-stored reading progress and JWT-based authentication makes
    the reading experience seamlessly resumable on any device on
    which the reader signs in.
    \item \textbf{Offline-friendly QR confirmation.} The librarian's
    confirmation flow requires only a smartphone with a camera; no
    additional hardware (barcode reader, magnetic-stripe terminal)
    is presupposed.
    \item \textbf{AI-assisted reading.} Selection-driven prompts allow
    on-the-fly explanation and translation, particularly relevant
    for readers approaching an unfamiliar language or a
    domain-specific text.
    \item \textbf{Reproducible deployment.} The Docker Compose stack
    is self-contained: \texttt{docker compose up -d --build} on a
    clean host is sufficient to bring the system online once the
    environment file is populated.
\end{itemize}

\section{Use Cases and Scenarios}

\textbf{Scenario 1 --- Student reserving a textbook.} A student
searches the catalogue for a textbook, narrows results by language and
genre, opens the detail page, observes that the campus branch holds
two available copies, and submits a reservation. The system mints a QR
code; the student visits the branch the next morning, the librarian
scans the code, and the reservation moves to \textit{active} with a
due date 14 days out. A \textit{return\_deadline} notification is
delivered two days before the due date.

\textbf{Scenario 2 --- Casual reader with a digital-only preference.}
A casual reader signs in, opens the catalogue, filters to
digital-available titles, and selects a book in EPUB. They read in
dark mode for several sessions; progress is synchronised on each
chapter boundary. When they encounter an unfamiliar passage in
Kazakh, they highlight it and ask the AI to translate to Russian. The
conversation is preserved in the AI sidebar for later reference.

\textbf{Scenario 3 --- Librarian processing returns.} The librarian
opens the scanner view, scans a returned book's QR code, and observes
the reservation transition to \textit{completed}. The
available-copies counter is restored automatically. A summary view
lists all today's \textit{active} reservations and their due dates.

\textbf{Scenario 4 --- Cross-device reading continuity.} A reader
begins a long-form Kazakh-language novel on their laptop in the
evening, advances to roughly 23\% of the book, and closes the tab.
The next morning, on a phone, signing in immediately resumes the
session at 23\% with the same chapter and the same highlights, with
no manual bookmark management required.

\section{Evaluation of Results}

The system was evaluated through (i) functional walk-throughs that
exercised each user story end-to-end, (ii) unit tests of pure utility
modules and form-validation schemas using Vitest, (iii) Go unit tests
of the JWT middleware and the SQL repositories, and (iv) exploratory
testing of the reader on representative EPUB and PDF files of varying
sizes. The functional walk-throughs confirmed that all
\textit{must}-tier stories from the backlog --- registration, login,
catalogue browsing, book detail, reservation, QR confirmation, EPUB
and PDF reading with progress persistence, wishlist, reviews, AI
conversation, and notification delivery --- complete without errors.

Performance was measured informally on a developer laptop and on the
production VM. On a representative book-detail request the end-to-end
latency, including database round-trip and JSON serialisation,
remained below 150 ms; the reader's first contentful render of an
EPUB chapter completed within roughly one second on a warm cache;
the AI streaming endpoint produced the first token within roughly
700 ms of the request, with the rest of the response delivered as
the upstream model emitted it.

\textit{[FIGURE PLACEHOLDER 4.15: Lighthouse audit screenshot of the
production deployment, captured for the catalogue page. Insert the
four score gauges (Performance, Accessibility, Best Practices, SEO).]}

\section{Justification of the Effectiveness of the Solution}

A direct comparison against the competitive baseline established in
Chapter 2 confirms the integration claim. None of the surveyed
platforms simultaneously offers a unified reader account, in-browser
EPUB and PDF reading, physical-copy reservation, and AI-assisted
reading in a Kazakh-first interface; Oqyrman delivers all four. The
QR-based confirmation workflow eliminates the operational friction of
manual reservation confirmation observed in Koha and ABIS deployments;
the percentage-based progress model preserves resumability across
reflow that page-number-based readers cannot offer; the
single-flight refresh design eliminates the duplicate-refresh race
condition that plagues naive token clients.

The platform's effectiveness is further amplified by reusability: the
relational schema, the REST contract, and the reader components are
positioned to be reused or extended without revising the
authentication or reservation subsystems, which positions the
platform as a foundation for further work rather than a one-off
deliverable.

\section{Discussion}

Several observations emerged during the iterative refinement of the
system that are worth recording for subsequent work.

\textit{Latency dependence on the upstream LLM.} The streaming UX is
functional but its perceived performance is dominated by upstream
latency. A retry policy with exponential backoff for transient
errors, plus a client-side timeout that surfaces a graceful
``the assistant is taking longer than usual'' state, would harden
the experience.

\textit{Percentage-based progress versus page numbers.} The
percentage-and-CFI model is robust against reflow but unfamiliar to
readers accustomed to ``page \textit{x} of \textit{y}'' presentation.
The reader currently displays both the percentage and the chapter
title, which testers found acceptable but unfamiliar at first
encounter.

\textit{Onboarding non-technical librarians.} Although the
librarian-side scanner is intentionally simple, librarians without
prior smartphone-camera-based workflows benefited from a brief
walk-through. A short embedded tutorial on first sign-in would close
this gap without increasing operational complexity.

\section{Risks and Mitigations}

Several risks were identified during planning and reviewed at each
sprint retrospective. Table~\ref{tab:risks} catalogues the principal
risks and the mitigations adopted.

\begin{table}[H]
\centering
\caption{Risk register and mitigations.}
\label{tab:risks}
\small
\begin{tabularx}{\textwidth}{l X X}
\toprule
\textbf{Risk} & \textbf{Description} & \textbf{Mitigation} \\
\midrule
Authentication compromise &
Stolen tokens or replayed refresh tokens granting indefinite access. &
Short access-token TTL (15 min); server-side refresh-token allowlist
with single-use semantics; bcrypt password hashing. \\
\addlinespace
LLM provider outage &
The upstream chat-completion endpoint is unavailable or degraded. &
Server-side proxy isolates the client from provider-specific
failures; retry with exponential backoff; fail soft to a typed error
that the UI surfaces gracefully. \\
\addlinespace
Database lock contention on hot rows &
\texttt{available\_copies} on popular titles becoming a hot row under
reservation load. &
All writes occur inside short transactions with row-level locks; the
counter sits on a single integer column with no triggers. \\
\addlinespace
Inventory drift between API and physical reality &
A book is physically returned without scanning. &
Periodic reconciliation job comparing
\texttt{available\_copies} against the physical count; visible
``mark-as-returned'' action for librarians. \\
\addlinespace
Catalogue translation drift &
The Kazakh and Russian descriptions diverge in meaning. &
Translations are stored side-by-side in the same row, reviewed
together; future work could connect a translation-memory tool. \\
\addlinespace
Adversarial prompts to the AI &
A user attempts to use the assistant for tasks unrelated to reading. &
A pinned system message constrains the assistant; prompt-length
limits cap the attack surface; logs preserve correlation
identifiers for after-the-fact review. \\
\bottomrule
\end{tabularx}
\end{table}

\section{Recommendations for Further Development}

Several directions extend naturally from the present work and would
substantially increase the value the platform delivers.

\begin{itemize}[leftmargin=*]
    \item \textbf{Audiobook support.} A third reader adapter alongside
    EPUB and PDF, with progress tracked in seconds rather than
    percentage, would broaden the catalogue to include narrated
    Kazakh classics and academic lectures.
    \item \textbf{Native mobile clients.} A React-Native client
    sharing the existing API would close the gap on commute and
    bedside reading scenarios where a browser-based reader is less
    convenient. Because the API is already mobile-friendly, the
    incremental cost is bounded by the client implementation.
    \item \textbf{Retrieval-augmented AI.} Indexing the user's
    library of read books and grounding assistant answers in that
    corpus would make the AI subsystem qualitatively more useful for
    academic readers, who increasingly expect citation-grounded
    answers.
    \item \textbf{Librarian analytics dashboard.} Reservation
    throughput, overdue ratios, popularity over time, and
    average-time-to-pickup are at present queryable only ad-hoc; a
    dedicated dashboard backed by aggregate queries against the
    existing tables would convert the platform's data into
    actionable acquisition guidance.
    \item \textbf{Push notifications.} The in-app notification inbox
    already exists; layering Web Push (and, for native clients, FCM)
    on top of the existing notification table is an additive
    extension that significantly improves time-to-awareness for
    deadlines.
    \item \textbf{Offline reading mode.} For readers in regions with
    intermittent connectivity, a service-worker-backed offline
    cache of currently-open books would close the last gap to
    parity with native readers.
\end{itemize}

% =============================================================================
%  CHAPTER 5 — CONCLUSION
% =============================================================================
\chapter{Conclusion}

\section{Summary of Achievements}

The thesis set out to deliver an integrated digital library platform
that unifies catalogue browsing, physical reservation, in-browser
reading, and AI-assisted reading support under a single multilingual
reader account. All eight objectives stated in the introduction were
met. A comparative study of international and domestic platforms
identified the integration gap that motivated the work; a normalised
PostgreSQL schema was designed and materialised through two canonical
migration files; a Go/Gin REST API was implemented with JWT
authentication, single-flight refresh-token rotation, and role-based
access control for librarian-only operations; a React 18 single-page
application was built on top of the API, exposing the catalogue, the
reservation workflow, the EPUB and PDF reader, the wishlist, the
reviews surface, and the notification inbox in three languages; an
LLM-backed assistant was integrated both as a global widget and as a
selection-driven layer inside the reader; a QR-based confirmation
workflow was implemented end-to-end; security was treated as a
cross-cutting concern aligned with the OWASP Top Ten; and the entire
system was containerised and put behind a continuous-delivery
pipeline that redeploys production on every push to \texttt{main}.

\section{Practical Impact}

The platform is directly applicable to the operational context of a
Kazakh public or academic library. Because the inventory model is
per-branch, a multi-branch library can be onboarded by populating
the \texttt{libraries} and \texttt{library\_books} tables without any
code change. Because the reader is format-driven, additional digital
editions can be attached to existing bibliographic records without
disturbing the rest of the schema. Because authentication is
stateless and the deployment is containerised, the system can be
scaled horizontally by running additional API replicas behind the
existing nginx; PostgreSQL becomes the single point at which
vertical scaling considerations apply, which is the conventional
locus for this class of application. The AI subsystem is portable
across LLM providers, since the server adapter speaks an
OpenAI-compatible chat contract.

From a societal standpoint, the platform contributes to the digital
inclusion of Kazakh-language readers by treating Kazakh as a
first-class interface language across every reader-facing surface,
and by unifying physical and digital reading under one identity in a
way that no surveyed competitor currently achieves.

\section{Closing Remarks}

In its current form Oqyrman already constitutes a coherent,
deployable answer to the integration gap identified in the
literature review, and provides a foundation on which the extensions
enumerated in Section 4.7 can be built incrementally. The codebase,
the schema, and the deployment pipeline are positioned to outlive the
diploma project itself, either as a continuing open-source effort or
as the technical core of a partnership with a domestic library
network. In an era when the physical and the digital sides of
literary culture are increasingly perceived as a single experience
by the reader, platforms that respect that perception --- and that
present three Kazakhstani languages with parity --- are not merely
convenient products but acts of cultural infrastructure.

% =============================================================================
%  REFERENCES
% =============================================================================
\begin{thebibliography}{99}
\addcontentsline{toc}{chapter}{References}

\bibitem{unesco2003}
UNESCO. \textit{Charter on the Preservation of the Digital Heritage}.
Adopted at the 32nd session of the General Conference of UNESCO,
2003. Available:
\url{https://unesdoc.unesco.org/ark:/48223/pf0000179529}

\bibitem{owasp2021}
OWASP Foundation. \textit{OWASP Top 10:2021}. The Open Web Application
Security Project, 2021. Available: \url{https://owasp.org/Top10/}

\bibitem{rfc7519}
M.~Jones, J.~Bradley, and N.~Sakimura. ``JSON Web Token (JWT).''
\textit{RFC 7519}, Internet Engineering Task Force, May 2015.
Available: \url{https://www.rfc-editor.org/rfc/rfc7519}

\bibitem{nielsen1994}
J.~Nielsen. \textit{Usability Engineering}. Morgan Kaufmann, 1994.

\bibitem{christen2012}
K.~Christen. ``Does Information Really Want to be Free? Indigenous
Knowledge Systems and the Question of Openness.''
\textit{International Journal of Communication}, vol.~6, pp.~2870--2893,
2012.

\bibitem{europeana2020}
Europeana Foundation. \textit{Europeana Strategy 2020--2025: A
Common Culture}. The Hague, 2020. Available:
\url{https://pro.europeana.eu/}

\bibitem{vesselinov2012}
R.~Vesselinov and J.~Grego. ``Duolingo Effectiveness Study.''
City University of New York, 2012.

\bibitem{licorish2018}
S.~A.~Licorish, H.~E.~Owen, B.~Daniel, and J.~L.~George. ``Students'
perception of Kahoot!'s influence on teaching and learning.''
\textit{Research and Practice in Technology Enhanced Learning},
vol.~13, no.~1, pp.~1--23, 2018.

\bibitem{idpf2017}
International Digital Publishing Forum. \textit{EPUB 3.2
Specification}. W3C, 2019. Available:
\url{https://www.w3.org/publishing/epub32/}

\bibitem{fielding2000}
R.~T.~Fielding. ``Architectural Styles and the Design of
Network-based Software Architectures.'' Ph.D.\ dissertation,
University of California, Irvine, 2000.

\bibitem{gohugoauth}
A.~Donovan and B.~W.~Kernighan. \textit{The Go Programming Language}.
Addison-Wesley Professional, 2015.

\bibitem{reactdocs}
Meta Open Source. \textit{React 18 Documentation}. 2023. Available:
\url{https://react.dev/}

\bibitem{tanstackquery}
TanStack. \textit{TanStack Query v5 Documentation}. 2024. Available:
\url{https://tanstack.com/query/latest}

\bibitem{postgresql}
The PostgreSQL Global Development Group. \textit{PostgreSQL 16
Documentation}. 2024. Available:
\url{https://www.postgresql.org/docs/16/}

\bibitem{dockerdocs}
Docker, Inc. \textit{Docker Compose Specification}. 2024. Available:
\url{https://docs.docker.com/compose/}

\end{thebibliography}

\end{document}
